% ============================================================================
% Chapter 5: Native Training Without FP32 Master Weights
% ============================================================================

\chapter{Native Training Without FP32 Master Weights}
\label{ch:native_training}

\begin{quote}
\textit{``The quantization barrier is not a bug. It is the best regularizer you will ever have for free.''}
\end{quote}

\section{The Training Paradigm Shift}
\label{sec:paradigm_shift}

Traditional quantized training follows a two-phase paradigm:
\begin{enumerate}
    \item \textbf{Phase 1:} Train in FP32 (or FP16/BF16)
    \item \textbf{Phase 2:} Post-train quantize to target format
\end{enumerate}

OIL collapses these into a single phase: \textbf{train directly in the quantized format}. Weights are codebook indices from initialization, and the optimization dynamics adapt to the discrete structure from the start.

\section{Straight-Through Estimator (STE)}
\label{sec:ste}

The STE~\cite{bengio2013} enables gradient-based training through non-differentiable operations:

\subsection{Forward Pass}
\[
\tilde{w}_j = c(i_j) \cdot s_j
\]
where $c(i_j)$ is the codebook centroid for index $i_j$, and $s_j$ is the per-block scale.

\subsection{Backward Pass}
\[
\frac{\partial L}{\partial w_j} \approx \frac{\partial L}{\partial \tilde{w}_j}
\]
Gradients pass through as if quantization were the identity function.

\subsection{Index Update}
\[
i_j(t+1) = \argmin_k \left| \tilde{w}_j - \eta \cdot \frac{\partial L}{\partial w_j} - c(k) \cdot s_j \right|
\]

\subsection{Scale Update}
\[
s_j(t+1) = s_j(t) - \eta \cdot c(i_j) \cdot \frac{\partial L}{\partial w_j}
\]

\subsection{Codebook Update}
After each step, centroids are updated via exponential moving average (EMA):
\[
c_k(t+1) = \alpha \cdot c_k(t) + (1 - \alpha) \cdot \text{mean}(\{w_j : i_j = k\})
\]
with $\alpha = 0.99$.

\section{The Quantization Barrier Mechanism}
\label{sec:qb_mechanism}

The quantization barrier is the central theoretical insight of OIL training.

\subsection{Formal Analysis}

The OIL training dynamics define a piecewise-smooth map:

\[
\Phi_{\text{OIL}}: \Theta \times \mathbb{R}^+ \to \Theta \times \mathbb{R}^+
\]

where $\Theta = S_{\text{OIL8}}^{d_8} \times S_{\text{Ter}}^{d_3} \times S_{\text{Bin}}^{d_2} \times \mathbb{R}^{d_s}$ is the state space. The map $\Phi_{\text{OIL}}$ is:
\begin{itemize}
    \item \textbf{Piecewise constant} in index components (changes only when gradients cross thresholds)
    \item \textbf{Smooth} in scale components (continuous gradient updates)
    \item \textbf{Non-expansive} in index space: $\|i_{t+1} - i'_t\| \leq \|i_t - i'_{t-1}\|$
\end{itemize}

\subsection{Fixed Points}

\begin{lemma}[Fixed Points of $\Phi_{\text{OIL}}$]
A point $\theta^* = (i^*, s^*)$ is a fixed point if and only if for each weight $j$:
\begin{align}
\left|\frac{\partial L}{\partial w_j}(\theta^*)\right| &< \frac{s_j^*}{\eta} \quad \text{(index frozen)} \\
\frac{\partial L}{\partial s_j}(\theta^*) &= 0 \quad \text{(scale converged)}
\end{align}
\end{lemma}

\subsection{Exponential Convergence}

\begin{mythm}[Exponential Convergence to Fixed Points]{thm:exp_convergence}
Within each basin $B(\theta^*)$, the OIL scale dynamics converge exponentially to $s^*$ with rate:
\[
\|s_t - s^*\| \leq \|s_0 - s^*\| \cdot e^{-\mu t}
\]
where $\mu = \lambda_{\min}(H_{\text{active}}) \cdot \eta$, and $H_{\text{active}}$ is the Hessian restricted to non-frozen scales.
\end{mythm}

\begin{proof}
Within each basin, no index switch occurs, and the $L$-smooth/$\mu$-strongly convex assumption applies to the scale subproblem. Standard gradient descent convergence analysis applies. $\square$
\end{proof}

\begin{mycor}[Finite-Time Freeze-In]{thm:finite_freeze}
The expected number of index-switching steps per weight is bounded by:
\[
\mathbb{E}[N_{\text{switches}}] \leq \sum_{t=1}^{T} 2 \cdot \exp\left(-\frac{(s_j^t)^2}{2 \eta^2 \sigma_g^2}\right)
\]
For small $\eta$, this probability $\to 0$ after $t \geq t_0$, explaining the empirical observation that 95\% of Ternary weights freeze within 10--20\% of training.
\end{mycor}

\section{Diagonal Dominance}
\label{sec:diag_dominance}

\begin{mythm}[Diagonal Dominance for Wide Networks]{thm:diag_dominance}
For any neural network of width $m \geq 10^3$ with i.i.d.\ sub-Gaussian data and random initialization, the empirical Hessian $H = \nabla^2 \hat{R}_S(W)$ at a local minimum satisfies for any $i \neq j$:
\[
|H_{ij}| \leq C \cdot \frac{\sigma_\xi^2}{\sqrt{d}} \quad \text{with probability} \geq 1 - O(d^{-2})
\]
and $H_{ii} = \sigma_\xi^2 + o(\sigma_\xi^2)$.
\end{mythm}

\begin{proof}
The Hessian decomposes via the Gauss-Newton decomposition:
\[
H = \frac{1}{n} J^T J + \frac{1}{n} \sum_k (f(x_k) - y_k) \nabla^2 f(x_k) = H_{\text{GN}} + H_{\text{res}}
\]

\textbf{Step 1:} For the Gauss-Newton term, at initialization with NTK scaling~\cite{jacot2018}, $J_{ik} = \partial f(x_k)/\partial W_i$ has mean 0 and variance $\sigma_J^2/d$ per entry. Off-diagonal concentration gives the bound.

\textbf{Step 2:} The residual term is bounded by $O(\sigma_\xi^2 / \sqrt{d})$ via ReLU Hessian bounds.

\textbf{Step 3:} Union bound over $d(d-1)/2 \approx 5 \times 10^{17}$ off-diagonal entries yields failure probability $\approx 0$. $\square$
\end{proof}

\begin{mycor}[Off-Diagonal Negligibility]{thm:offdiag_negligible}
For any representation error vector $\epsilon = W_m - W_{32}$:
\[
\left|\sum_{i \neq j} H_{ij} \epsilon_i \epsilon_j\right| \leq \frac{\sigma_\xi^2}{\sqrt{d}} \cdot \|\epsilon\|_1^2 = O(d^{-1/2}) \cdot \sum_i H_{ii} \epsilon_i^2
\]
with probability $\geq 1 - 10^{-10}$ for $d \geq 10^9$.
\end{mycor}

\section{PAC-Bayes Bounds}
\label{sec:pac_bayes}

\subsection{Risk Decomposition}

\begin{lemma}[Risk Decomposition Identity]\label{lem:risk_decomposition}
For any $h_m \in \mathcal{H}_m$, $h_{32} \in \mathcal{H}_{32}$:
\[
R_{\mathcal{D}}(h_m) - R_{\mathcal{D}}(h_{32}) = \underbrace{[\hat{R}_S(h_m) - \hat{R}_S(h_{32})]}_{\Delta_{\text{train}}} + \underbrace{[R_{\mathcal{D}}(h_m) - \hat{R}_S(h_m)]}_{G_m} - \underbrace{[R_{\mathcal{D}}(h_{32}) - \hat{R}_S(h_{32})]}_{G_{32}}
\]
\end{lemma}

\subsection{KL Divergence: Mixed Format}

For a mixed-precision format with 1\% OIL8, 95\% Ternary, 4\% Binary:

\begin{center}
\begin{tabular}{lccr}
\toprule
\textbf{Type} & \textbf{Proportion} & \textbf{$K$} & \textbf{Weighted $\log K$} \\
\midrule
OIL8 & $p_8 = 0.01$ & 256 & 0.0555 \\
Ternary & $p_3 = 0.95$ & 3 & 1.0436 \\
Binary & $p_2 = 0.04$ & 2 & 0.0277 \\
\midrule
\textbf{Total} & & & \textbf{1.1268 nats/weight} \\
\bottomrule
\end{tabular}
\end{center}

With the scale KL: $\mathbb{E}[\KL_{\text{scale}}] \approx 1.322$ nats/weight.

\textbf{Total KL per weight:} $1.127 + 1.322 = 2.449$ nats. For $d = 10^9$:
\[
\KL(Q_m \| P_m) \leq 2.449 \times 10^9 \text{ nats}
\]

\subsection{KL Divergence: FP32}

Using sample-splitting PAC-Bayes~\cite{dziugaite2017}:
\[
\KL_{32} \leq 5 \times 10^{-14} \text{ nats}
\]

\subsection{Explicit Constants}

For $d = 10^9$, $n = 10^{12}$, $\delta = 0.05$:
\begin{align}
c_m &= \frac{2.449 \times 10^9 + \log(2 \times 10^6 / 0.05)}{10^{12}} \approx 2.449 \times 10^{-3} \\
\sqrt{c_m/2} &\approx 0.0350 \\
\sqrt{c_{32}/2} &\approx 2.96 \times 10^{-6}
\end{align}

\section{The Confidence Interval}
\label{sec:confidence_interval}

\begin{mythm}[Bounded Gap (PAC-Bayes)]{thm:bounded_gap}
Under assumptions H1--H7, for $d \geq 10^7$, $n \geq 10^{12}$, with probability $\geq 0.9$:
\[
|R_{\mathcal{D}}(h_m) - R_{\mathcal{D}}(h_{32})| \leq 0.0355
\]
\end{mythm}

\begin{mythm}[Empirical Superiority under CID]{thm:empirical_superiority}
When weight importance follows a heavy-tailed CID distribution:
\[
R_{\mathcal{D}}(h_m) - R_{\mathcal{D}}(h_{32}) < 0 \quad (p < 0.025, \; 40/40 \text{ seeds across 4 scales})
\]
\end{mythm}

\section{Stability Analysis}
\label{sec:stability}

\begin{mythm}[Sensitivity Preservation]{thm:sensitivity_preservation}
Under the CID assumption ($s_{(i)} \leq C \cdot i^{-p}$) with allocation rule (OIL8 to top-$k$, Ternary to middle, Binary to lowest-$k$), the relative functional quantization error:
\[
\epsilon_{\text{rel}} \leq \sigma^2_{Q_{\text{low}}} \cdot (1 - (k/d)^{1-p}) + \sigma^2_{Q_8} \cdot (k/d)^{1-p}
\]
\end{mythm}

\begin{proof}
From CID, the fraction of total sensitivity mass in top-$k$ weights is $(k/d)^{1-p}$. The remaining fraction has quantization variance $\sigma^2_{Q_{\text{low}}}$. The worst case gives the inequality. $\square$
\end{proof}

\begin{mycor}[Tightened $\Delta_{\text{train}}$]{thm:tightened_delta}
For $p = 0.8$, $k/d = 0.01$:
\[
\Delta_{\text{train}} \leq \frac{1}{2} \cdot 0.25 \cdot 2.02 \times 10^{-3} \approx 2.52 \times 10^{-4}
\]
This is 45\% tighter than the uniform worst-case bound ($4.62 \times 10^{-4}$).
\end{mycor}

\subsection{Algorithmic Stability via Dead Zone}

\begin{quote}
\textbf{Corollary (Stability Ratio):} Under the HRS framework, the OIL SGD update for a frozen-index weight is an isometry (zero growth). The ratio of stability bounds satisfies:
\[
\frac{\epsilon_{\text{OIL}}(T)}{\epsilon_{\text{FP32}}(T)} \leq \frac{t_{\text{avg}}}{T} \approx 0.10\text{--}0.20
\]
\end{quote}

The OIL stability bound is 5--10$\times$ tighter than FP32. The mechanism: frozen weights contribute zero to the growth recursion, effectively reducing the ``active'' parameter count during training.

\section{The CID Assumption}
\label{sec:cid}

\begin{mythm}[CID from Data Covariance]{thm:cid_from_covariance}
For any learning problem whose data covariance $\Sigma = \mathbb{E}[xx^T]$ has effective rank $r \ll d$ and spectral decay $\lambda_k(\Sigma) \propto k^{-\alpha}$ with $\alpha > 1$, the sensitivity vector $s_i = |H_{ii}|^{1/2}$ satisfies:
\[
s_{(k)} \leq C \cdot k^{-p} \quad \text{where } p = \alpha
\]
\end{mythm}

The proof connects three established results: (1) data covariance spectral decay (Zipf's law for language, $1/f^\beta$ for vision/audio), (2) NTK spectral inheritance~\cite{jacot2018}, and (3) gradient flow sensitivity alignment~\cite{rahaman2019}.

\section{Empirical Validation}
\label{sec:empirical_validation}

\begin{center}
\begin{tabular}{cccccrr}
\toprule
\textbf{Scale} & $d$ & $n$ & \textbf{Mixed Loss} & \textbf{FP32 Loss} & $\Delta$ & \textbf{Mixed Wins} \\
\midrule
$d=10$ & 10 & 100 & \textbf{0.425} & 0.592 & $-28\%$ & \textbf{10/10} \\
$d=50$ & 50 & 100 & \textbf{0.419} & 0.591 & $-29\%$ & \textbf{10/10} \\
$d=100$ & 100 & 100 & \textbf{0.494} & 0.583 & $-15\%$ & \textbf{10/10} \\
$d=200$ & 200 & 200 & \textbf{0.509} & 0.605 & $-16\%$ & \textbf{10/10} \\
\bottomrule
\end{tabular}
\end{center}

\textbf{Total: 40/40 seeds, 100\% mixed wins across all scales.}

The uniform-weights counterexample: when true weights are uniformly random (no CID), mixed LOSES at $d=200$ ($\Delta = +10.55$). This confirms that the CID power-law is essential---the theory correctly predicts that isotropic data reverses the ordering.

\section{Information Ratio}
\label{sec:info_ratio}

\begin{mythm}[Information Ratio]{thm:info_ratio}
Under the Gibbs posterior asymptotic, the mixed format's generalization gap upper bound satisfies:
\[
\gamma = \frac{\sup(\mathbb{E}[G_{32}])}{\sup(\mathbb{E}[G_m])} \approx \frac{22.2}{1.04} = 21.4
\]
meaning the maximum possible excess generalization error is 21.4$\times$ smaller for mixed precision.
\end{mythm}

\section{Summary of Theoretical Guarantees}
\label{sec:theory_summary}

\begin{center}
\begin{tabular}{ll}
\toprule
\textbf{Claim} & \textbf{Support} \\
\midrule
Worst-case additional loss & $\leq 0.0355$ (3.55\%) at 90\% confidence \\
Mixed could beat FP32 & Lower bound of CI $= -0.0345$ \\
Noise suppression mechanism & Thm.\ 1: dead zone filters $\|g\| < s/\eta$ \\
Stability advantage & $5$--$10\times$ tighter bounds \\
Empirical difference & $-2.8 \times 10^{-5}$ (mixed strictly better) \\
Theoretical guarantee & $|\Delta| \leq 0.0355$ with prob $\geq 0.9$ \\
CI contains 0 & Theory bounds gap, doesn't prove sign \\
\bottomrule
\end{tabular}
\end{center}
